% Validación del modelo de negocio, entrega final. Máximo 3 páginas, sin contar
% el anexo.
% Sustituir cada texto gris (\pendiente{...}) por el contenido del equipo.
% No cambiar los títulos ni el orden de las secciones.
% Las cifras, las tablas y el anexo se generan solos a partir de hipotesis.csv,
% entrevistas.csv y metricas_landing.csv. No escribir números a mano: actualizar
% los CSV y volver a compilar.
% Compilar dos veces con: pdflatex validacion.tex
\documentclass[11pt]{article}
\usepackage[utf8]{inputenc}
\usepackage[T1]{fontenc}
\usepackage[spanish,es-noindentfirst,es-nodecimaldot,es-tabla]{babel}
\usepackage[letterpaper,margin=2.2cm,top=1.6cm]{geometry}
\usepackage{xcolor}
\usepackage{booktabs,array,titling,enumitem,float,longtable,datatool}
\usepackage[font=small,labelfont=bf,skip=4pt]{caption}
\usepackage[hidelinks]{hyperref}

\definecolor{unamAzul}{HTML}{003D79}
\definecolor{unamOro}{HTML}{D59F0F}
\newcommand{\pendiente}[1]{\textcolor{gray}{\emph{#1}}}
\newcolumntype{L}[1]{>{\raggedright\arraybackslash}p{#1}}
\setlist{nosep,leftmargin=*}
\setlength{\parskip}{4pt}
\setlength{\parindent}{0pt}

% ---------- Datos ----------
% csv-content=literal permite %, &, $ y _ en los textos de los CSV. Las claves
% siguen el orden de columnas de la plantilla, que no se cambia.
\DTLread[name=hip,format=csv,csv-content=literal,
  keys={id,bloque,hipotesis,prueba,criterio,estado,evidencia,cambio}]{hipotesis.csv}
\DTLread[name=ent,format=csv,csv-content=literal,
  keys={id,fecha,perfil,sector,canal,duracion,hipotesis,dolor,pagahoy,pagaria,referidos,
        hallazgo,cita}]{entrevistas.csv}
\DTLread[name=lan,format=csv,csv-content=literal,
  keys={semana,visitas,registros,fuente,cambio}]{metricas_landing.csv}

% Indicadores de las entrevistas
\newcounter{nent}\newcounter{ndolor}\newcounter{npaga}\newcounter{npagaria}
\newcounter{nval}\newcounter{ninval}\newcounter{npend}
% Indicadores de la landing page
\newcounter{nsem}\def\totVis{0}\def\totReg{0}

% ---------- Encabezado institucional ----------
\pretitle{%
  \begin{center}\small
    Universidad Nacional Autónoma de México\\
    Posgrado en \pendiente{nombre del programa}\\[2pt]
    \pendiente{Nombre del curso}, semestre \pendiente{2027-1}
  \end{center}
  \vspace{-4pt}{\color{unamOro}\rule{\textwidth}{1.2pt}}\par\vspace{14pt}
  \begin{center}{\small\scshape Validación del modelo de negocio}\\[6pt]\LARGE\bfseries\color{unamAzul}}
\posttitle{\par\end{center}\vspace{-4pt}}
\preauthor{\begin{center}\large\begin{tabular}[t]{c}}
\newcommand{\otrorenglon}{\end{tabular}\par\vspace{8pt}\begin{tabular}[t]{c}}
\postauthor{\end{tabular}\par\end{center}}
\predate{\begin{center}\small}
\postdate{\par\end{center}}

\title{\pendiente{Nombre de la empresa: el proyecto en una frase}}
% Los mismos integrantes de la propuesta. Borrar los que sobren, junto con su \and.
\author{%
  \pendiente{Integrante 1}\\ {\small\texttt{\pendiente{correo}}}
  \and
  \pendiente{Integrante 2}\\ {\small\texttt{\pendiente{correo}}}
  \and
  \pendiente{Integrante 3}\\ {\small\texttt{\pendiente{correo}}}
}
\date{\today}

\begin{document}
\maketitle
\thispagestyle{empty}

\DTLforeach*{ent}{\eDolor=dolor,\ePaga=pagahoy,\ePagaria=pagaria}{%
  \stepcounter{nent}%
  \DTLifnumgt{\eDolor}{3}{\stepcounter{ndolor}}{}%
  \DTLifstringeq{\ePaga}{si}{\stepcounter{npaga}}{}%
  \DTLifstringeq{\ePagaria}{si}{\stepcounter{npagaria}}{}}
\DTLforeach*{hip}{\hEstado=estado}{%
  \DTLifstringeq{\hEstado}{validada}{\stepcounter{nval}}{}%
  \DTLifstringeq{\hEstado}{invalidada}{\stepcounter{ninval}}{}%
  \DTLifstringeq{\hEstado}{pendiente}{\stepcounter{npend}}{}}
\DTLforeach*{lan}{\lVis=visitas,\lReg=registros}{%
  \stepcounter{nsem}%
  \xdef\totVis{\fpeval{\totVis+\lVis}}\xdef\totReg{\fpeval{\totReg+\lReg}}}

% ---------------------------------------------------------------------------
% RESUMEN. Las cifras del primer párrafo se calculan solas. El segundo párrafo
% las interpreta.
% ---------------------------------------------------------------------------
\section{Resumen de la evidencia}
\label{sec:resumen}

El equipo hizo \thenent{} entrevistas con clientes potenciales. En \thendolor{} de ellas la
persona describió un dolor de 4 o 5 en una escala de 1 a 5, en \thenpaga{} ya gasta tiempo o
dinero en resolverlo y en \thenpagaria{} dio una señal concreta de compra. El
Anexo~\ref{sec:anexo} registra cada entrevista. La landing page, publicada en
% Sustituir por \url{https://...}
\pendiente{dirección de la landing page}, recibió \totVis{} visitas y \totReg{} registros en
\thensem{} semanas, una conversión de
\ifnum\totVis>0 \fpeval{round(100*\totReg/\totVis,1)}\else 0\fi\%. De las hipótesis del canvas,
\thenval{} quedaron validadas, \theninval{} invalidadas y \thenpend{} siguen pendientes.

\pendiente{Un párrafo que interpreta estos números: qué hipótesis quedaron validadas, cuál fue
la sorpresa más importante y qué falta probar.}

% ---------------------------------------------------------------------------
% HIPÓTESIS. Estado y evidencia de cada hipótesis de modelo_negocio.tex. E indica
% una entrevista y L una semana de la landing page.
% ---------------------------------------------------------------------------
\section{Resultado de las hipótesis}
\label{sec:hipotesis}

La Tabla~\ref{tab:hipotesis} muestra el resultado de cada hipótesis frente al criterio de
éxito fijado antes de probarla.

\begin{table}[H]
  \caption{Hipótesis del canvas, con su criterio de éxito, su estado y la evidencia que la
    respalda. E indica una entrevista y L una semana de la landing page.}
  \label{tab:hipotesis}
  \footnotesize
  \renewcommand{\arraystretch}{1.2}
  \begin{tabular}{@{}lL{0.38\textwidth}L{0.26\textwidth}L{0.09\textwidth}L{0.1\textwidth}@{}}
    \toprule
    id & Hipótesis (bloque) & Criterio de éxito & Estado & Evidencia \\
    \midrule
    \DTLforeach*{hip}{\hId=id,\hBloque=bloque,\hTexto=hipotesis,\hCrit=criterio,%
                      \hEstado=estado,\hEvid=evidencia}{%
      \DTLiffirstrow{}{\\}%
      \hId & \hTexto{} {\color{gray}(\hBloque)} & \hCrit & \hEstado &
      \DTLsubstituteall{\hEvid}{;}{; }\hEvid}\\
    \bottomrule
  \end{tabular}
\end{table}

% ---------------------------------------------------------------------------
% LANDING PAGE. La tabla se genera a partir de metricas_landing.csv.
% ---------------------------------------------------------------------------
\section{Landing page}
\label{sec:landing}

\pendiente{Una o dos oraciones que dicen qué cambio en la página produjo el mayor aumento en la
conversión.} La Tabla~\ref{tab:landing} muestra las métricas de cada semana.

\begin{table}[H]
  \caption{Métricas semanales de la landing page. La conversión es registros entre visitas.}
  \label{tab:landing}
  \footnotesize
  \renewcommand{\arraystretch}{1.2}
  \setlength{\tabcolsep}{4pt}
  \begin{tabular}{@{}lL{0.1\textwidth}rrrL{0.25\textwidth}L{0.25\textwidth}@{}}
    \toprule
    & Semana & Visitas & Registros & Conversión & Fuente de tráfico & Cambio en la página \\
    \midrule
    \DTLforeach*{lan}{\lSem=semana,\lVis=visitas,\lReg=registros,\lFuente=fuente,%
                      \lCambio=cambio}{%
      \DTLiffirstrow{}{\\}%
      L\DTLcurrentindex & \lSem & \lVis & \lReg &
      \ifnum\lVis>0 \fpeval{round(100*\lReg/\lVis,1)}\%\else --\fi & \lFuente & \lCambio}\\
    \bottomrule
  \end{tabular}
\end{table}

% ---------------------------------------------------------------------------
% CAMBIOS EN EL CANVAS. La tabla muestra las hipótesis con algo en la columna
% "cambio" de hipotesis.csv. El párrafo cuenta cómo cambió el modelo desde la
% primera entrega. Una hipótesis invalidada que cambió el canvas es un buen
% resultado.
% ---------------------------------------------------------------------------
\section{Cambios en el canvas}
\label{sec:cambios}

\pendiente{Un párrafo que cuenta cómo cambió el modelo desde la primera entrega y qué evidencia
provocó cada cambio. Si nada cambió, explicar por qué.} La Tabla~\ref{tab:cambios} lista los
cambios, y el canvas actualizado está en \texttt{modelo\_negocio.pdf}.

\begin{table}[H]
  \caption{Cambios en el canvas y la hipótesis que los provocó.}
  \label{tab:cambios}
  \footnotesize
  \renewcommand{\arraystretch}{1.2}
  \begin{tabular}{@{}lL{0.12\textwidth}L{0.78\textwidth}@{}}
    \toprule
    id & Estado & Cambio en el canvas \\
    \midrule
    \DTLforeach*[\not\equal{\hCambio}{}]{hip}{\hId=id,\hEstado=estado,\hCambio=cambio}{%
      \DTLiffirstrow{}{\\}%
      \hId & \hEstado & \hCambio}\\
    \bottomrule
  \end{tabular}
\end{table}

% ---------------------------------------------------------------------------
% ANEXO. Se genera a partir de entrevistas.csv y no cuenta para el límite de
% páginas.
% ---------------------------------------------------------------------------
\clearpage
\appendix
\section{Registro de entrevistas}
\label{sec:anexo}

\begin{footnotesize}
\renewcommand{\arraystretch}{1.2}
\setlength{\tabcolsep}{4pt}
\begin{longtable}{@{}lL{0.1\textwidth}L{0.22\textwidth}cccL{0.39\textwidth}@{}}
  \caption{Entrevistas con clientes potenciales. Dolor de 1 (ninguno) a 5 (urgente). ``Paga''
    indica si ya gasta tiempo o dinero en el problema, y ``pagaría'' si dio una señal concreta
    de compra.}\label{tab:entrevistas}\\
  \toprule
  id & Fecha & Perfil & Dolor & Paga & Pagaría & Hallazgo y cita \\
  \midrule
  \endhead
  \DTLforeach*{ent}{\eId=id,\eFecha=fecha,\ePerfil=perfil,\eDolor=dolor,\ePaga=pagahoy,%
                    \ePagaria=pagaria,\eHallazgo=hallazgo,\eCita=cita}{%
    \DTLiffirstrow{}{\\}%
    \eId & \eFecha & \ePerfil & \eDolor & \ePaga & \ePagaria &
    \eHallazgo\newline{\color{gray}``\eCita''}}\\
  \bottomrule
\end{longtable}
\end{footnotesize}

\end{document}
